% !TEX root = ./main.tex

% =====================================================
% PREÂMBULO — pacotes, configurações e ambientes
% =====================================================

% ---------- Codificação e idioma ----------
\usepackage[utf8]{inputenc}
\usepackage[T1]{fontenc}
\usepackage[brazilian]{babel}
\usepackage{lmodern}
%\usepackage[alf]{abntex2cite}

% ---------- Matemática ----------
\usepackage{amsmath}
\usepackage{amssymb}
\usepackage{amsthm}
\usepackage{mathtools}
\usepackage{bm}               % negrito em símbolos matemáticos

% ---------- Figuras e desenhos ----------
\usepackage{graphicx}
\usepackage{tikz}
\usetikzlibrary{arrows.meta, calc, positioning, decorations.markings, angles, quotes}
\usetikzlibrary{shapes.geometric}
\usetikzlibrary{shapes.symbols}
\usetikzlibrary{shadows}
\usepackage{pgfplots}
\usepackage{float}
\usepackage{subcaption}

% ---------- Layout de página ----------
\usepackage[a4paper, top=3cm, bottom=2cm, left=3cm, right=2cm, headheight=14.5pt]{geometry}
\usepackage{setspace}
\onehalfspacing
\usepackage{indentfirst}
\usepackage{fancyhdr}

% ---------- Tabelas ----------
\usepackage{booktabs}
\usepackage{multirow}

% ---------- Referências cruzadas e citações ----------
\usepackage[
  backend=biber,
  style=abnt,
  sorting=none
]{biblatex}
\addbibresource{referencias.bib}

\usepackage{hyperref}
\usepackage[capitalise]{cleveref}
\hypersetup{
  colorlinks=true,
  linkcolor=black,
  citecolor=black,
  urlcolor=blue,
  pdftitle={Geometria Riemanniana: Geodésicas e o Teorema de Hopf-Rinow},
  pdfauthor={João Paulo Chiari de Gasperi, Welinton Anderson Rocha, Roger Emanuel Moraes Pezzott}
}

% ---------- Cabeçalho/rodapé ----------
\pagestyle{fancy}
\fancyhf{}
\fancyhead[R]{\thepage}
\renewcommand{\headrulewidth}{0.4pt}

% =====================================================
% AMBIENTES MATEMÁTICOS
% Numerados por capítulo (ex.: Teorema 2.1, Lema 2.2, ...)
% Todos compartilham o mesmo contador para manter a ordem
% cronológica de aparição no texto.
% =====================================================
\theoremstyle{plain}

\newtheorem{theorem}{Teorema}[chapter]
\newtheorem{proposition}{Proposição}[chapter]
\newtheorem{lemma}{Lema}[chapter]
\newtheorem{corollary}{Corolário}[chapter]

\theoremstyle{definition}

\newtheorem{example}{Exemplo}[chapter]
\newtheorem{examples}[example]{Exemplos}
\newtheorem{definition}{Definição}[chapter]
\newtheorem{axiom}{Axioma}[chapter]

\theoremstyle{remark}

\newtheorem{observation}{Observação}[chapter]

% Fim de demonstração customizado (símbolo QED)
\renewcommand{\qedsymbol}{$\blacksquare$}

% Nome em minúsculo (uso em \cref) e maiúsculo (uso em \Cref)
\crefname{axiom}{axioma}{axiomas}
\Crefname{axiom}{Axioma}{Axiomas}

\crefname{theorem}{teorema}{teoremas}
\Crefname{theorem}{Teorema}{Teoremas}

\crefname{proposition}{proposição}{proposições}
\Crefname{proposition}{Proposição}{Proposições}

\crefname{lemma}{lema}{lemas}
\Crefname{lemma}{Lema}{Lemas}

\crefname{corollary}{corolário}{corolários}
\Crefname{corollary}{Corolário}{Corolários}

\crefname{example}{exemplo}{exemplos}
\Crefname{example}{Exemplo}{Exemplos}

\crefname{examples}{exemplos}{exemplos}
\Crefname{examples}{Exemplos}{Exemplos}

\crefname{definition}{definição}{definições}
\Crefname{definition}{Definição}{Definições}

\crefname{observation}{observação}{observações}
\Crefname{observation}{Observação}{Observações}



%% Ambiente de citação padrão abnt

\newenvironment{citacao}{%
  \begin{list}{}{%
      \small % Fonte menor (geralmente tamanho 10)
      \setlength{\leftmargin}{4cm} % Recuo de 4 cm da margem esquerda
      \setlength{\rightmargin}{0cm}
      \setlength{\itemindent}{0cm}
      \setlength{\labelwidth}{0cm}
      \setlength{\labelsep}{0cm}
      \setlength{\parsep}{0cm}
      \setstretch{1.0} % Espaçamento simples (requer pacote setspace se usado)
    }%
    \item[]%
          }{%
  \end{list}%
}
